\Article{30}
The theorem of the preceding article needs some modification in a special case that we cannot pass over in silence. We have tacitly supposed that the surface of the vessel has continuous curvature along the whole boundary $P$, so that at every point of this boundary there is a \textit{single} plane tangent to the vessel’s surface. If the continuity of the curvature is interrupted at an isolated point of the line $P$, whether there is a cusp there or an edge crosses the line $P$, it is easy to see that our conclusion is unaffected. The situation is different when the continuity of the curvature is interrupted along a finite part of the line $P$, that is, when the surface of the vessel has an edge along a finite part of the line $P$ or even along that whole line. At every point of such a part there will then be two planes
\SourceLeaf{74}tangent to the surface of the vessel, one belonging to the free part of that surface and the other to the part $T$. If we retain the symbol $i$ for the angle between the first plane and the plane tangent to the surface $U$, and denote by $k$ the angle between this tangent plane and the second plane, we shall no longer have $i+k=180^0$. The sum will be greater or smaller according as the edge is convex or concave. The element of the variation $\delta W$ for a virtual motion in which the fluid spills beyond the boundary $P$ is still expressed by
\[
\alpha\alpha\mathrm{d}P.\delta e.\sin(5,6).(\cos A-\cos i)
\]
while for a virtual motion in which the fluid recedes within the boundary $P$, the element of that variation will now be
\[
-\alpha\alpha\mathrm{d}P.\delta e.\sin(5,6).(\cos A+\cos k)
\]
To prevent a negative variation in the value of $W$, the value of $\cos A-\cos i$ must be nonnegative and that of $\cos A+\cos k$ nonpositive. That is, we must have
\[
\begin{aligned}
&\text{either}\quad i=A,\qquad\text{either}\quad i>A\\
&\text{and either}\quad k=180^0-A,\qquad\text{either}\quad k>180^0-A
\end{aligned}
\]
In the state of equilibrium, therefore, we cannot have $i+k<180^0$. Equivalently, \textit{in the state of equilibrium the boundary of the free surface $U$ of the fluid cannot lie along a finite stretch of a concave edge of the vessel’s surface.} Conversely, whenever part of that boundary coincides with a convex edge, the necessary and sufficient condition for equilibrium is that the angle between the planes tangent to the fluid and the vessel lie between $A$ and $A+a$ (inclusive) when measured outside the fluid, or between $180^0-A$ and $180^0-A+a$ when measured within the fluid. Here $180^0-a$ denotes the angle at any point between the two planes tangent to the vessel’s surface on either side of the edge, measured on the vessel side.

\Article{31}
The constants $\alpha\alpha$, $\beta\beta$, whose ratio determines the angle $A$, depend on the functions $f,F$ and can in a certain sense be regarded as measures of the strength of the molecular forces exerted by the particles of the fluid and of the vessel. If these functions are such that $fx,Fx$ are in a fixed ratio independent of the distance $x$, say as $n$ to $N$, we can evidently set
\SourceLeaf{75}$\alpha\alpha:\beta\beta=cn:CN$. In other words, the constants $\alpha\alpha$, $\beta\beta$ will be proportional to the attractions exerted at the same distance by two molecules equal in volume, one of the fluid and the other of the vessel. Now the angle $A$ is acute, equal to a right angle, obtuse, or equal to two right angles, according as $\beta\beta<\frac12\alpha\alpha$, $\beta\beta=\frac12\alpha\alpha$, $\beta\beta>\frac12\alpha\alpha$ but $<\alpha\alpha$, or $\beta\beta=\alpha\alpha$. Under this assumption, which, though unsubstantiated, is not implausible, we must say that the first case occurs whenever the mutual attraction of the parts of the fluid is greater than twice the attraction of the parts of the vessel on the fluid. The second occurs whenever the former attraction is twice the latter. The third occurs whenever the former is greater than the latter but less than twice it. Finally, the fourth occurs whenever both attractions are equal. Mercury in glass vessels is an example of the first case.

